\clearpage

\subsection{Диапазон наблюдений, активация и восстановление}

Таблица при нагрузочных возвратах показывает выбор retained-load. Числа порогов обозначают фактические значения WIP и число обучений. Диапазон обучения — разброс максимального наблюдавшегося WIP между наборами; $q/\max w$ считается внутри каждого набора. K/D обозначают известную/обнаруженную структуру. Диагностика не меняет правила активации.

\begin{center}\small\setlength{\tabcolsep}{3pt}\begin{tabular}{llp{.24\textwidth}rrr}\toprule

Process & N & Threshold:count & Train max & $q/\max w$ & Above\\\midrule

P1 & K & 6:11; 12:1 & 5--8 & 0.75--2.40 & 4/12\\

P1 & D & 6:12 & 5--8 & 0.75--1.20 & 4/12\\

P2 & K & 8:12 & 6--12 & 0.67--1.33 & 2/12\\

P2 & D & 8:12 & 6--12 & 0.67--1.33 & 2/12\\

P3 & K & 10:11; none:1 & 7--11 & 0.91--1.43 & 5/12\\

P3 & D & 10:11; none:1 & 7--11 & 0.91--1.43 & 5/12\\

P4 & K & 10:12 & 11--19 & 0.53--0.91 & 0/12\\

P4 & D & 10:12 & 11--19 & 0.53--0.91 & 0/12\\

P5 & K & 10:11; 20:1 & 12--18 & 0.56--1.18 & 1/12\\

P5 & D & 10:11; 20:1 & 12--18 & 0.56--1.18 & 1/12\\

P6 & K & 10:11; none:1 & 9--13 & 0.77--1.11 & 1/12\\

P6 & D & 10:12 & 9--13 & 0.77--1.11 & 1/12\\

\bottomrule\end{tabular}\end{center}

\begin{center}\small\setlength{\tabcolsep}{4pt}\begin{tabular}{llrrrrr}\toprule

Process & N & Activate & Nonsevere & Lead/$T$ & Event & Prevented\\\midrule

P1 & K & 0.717 & 0.371 & 5.661 & 0.365 & 0.326\\

P1 & D & 0.724 & 0.385 & 5.979 & 0.378 & 0.299\\

P2 & K & 0.736 & 0.542 & 5.666 & 0.325 & 0.223\\

P2 & D & 0.736 & 0.542 & 5.666 & 0.370 & 0.117\\

P3 & K & 0.628 & 0.052 & 4.649 & 0.651 & 0.023\\

P3 & D & 0.628 & 0.052 & 4.649 & 0.651 & 0.023\\

P4 & K & 0.828 & 0.533 & 8.362 & 0.490 & 0.223\\

P4 & D & 0.828 & 0.533 & 8.362 & 0.490 & 0.223\\

P5 & K & 0.833 & 0.640 & 9.790 & 0.354 & 0.330\\

P5 & D & 0.833 & 0.640 & 9.790 & 0.354 & 0.330\\

P6 & K & 0.667 & 0.203 & 6.346 & 0.401 & 0.382\\

P6 & D & 0.729 & 0.227 & 6.302 & 0.528 & 0.187\\

\bottomrule\end{tabular}\end{center}

Nonsevere — активация при отсутствии тяжёлого события в парной траектории без действия; это не обязательно вредное вмешательство. Время предупреждения измеряется относительно начала такого контрфактического события. Предотвращение условно по событию без действия. Восстановление ниже означает WIP не выше $2S$ в течение $10T$ после импульса с цензурированием на исходном горизонте $180T$. Сначала усредняются траектории каждого независимого обучения, затем сами обучения; учитываются все импульсные траектории, а не только кризисные.

\begin{center}\small\begin{tabular}{lrrrr}\toprule

Process & No action & Predictive K & Retained K & Generator\\\midrule

P1 & 0.792 & 1.000 & 1.000 & 1.000\\

P2 & 0.742 & 0.958 & 1.000 & 1.000\\

P3 & 0.169 & 0.667 & 0.560 & 0.727\\

P4 & 0.555 & 0.979 & 0.878 & 0.992\\

P5 & 0.815 & 0.917 & 0.987 & 1.000\\

P6 & 0.477 & 0.935 & 0.953 & 1.000\\

\bottomrule\end{tabular}\end{center}

\subsection{Последствия общих ресурсов и положение узкого места}

Ресурсная диагностика повторяет P4/P5 при всех четырёх механизмах и всех 12 обучениях для двух импульсных условий с высокой базовой интенсивностью и повтора 0. Включены все типы действий на минимальном фиксированном пороге и сохранённые выбранные политики. Все 880 префиксов до $180T$ точно совпадают с исходными тестовыми отсчётами. Затем поступления прекращаются, а действия сохраняются до $600T$ для завершения когорты. Интегралы очередей относятся только к исходному горизонту. Поэтому длительность полной когорты имеет другой горизонт, чем основной показатель.

\begin{longtable}{llrrrr}\caption{P4: интегралы WIP каждого процесса (тысячи) по полному тесту и времена завершения полной когорты по продлённой диагностике.}\\\toprule

Mechanism & Action at 10 & WIP X & WIP Y & Cycle X & Cycle Y\\\midrule\endfirsthead

\toprule Mechanism & Action at 10 & WIP X & WIP Y & Cycle X & Cycle Y\\\midrule\endhead

fixed & None & 7.78 & 8.61 & 21.53 & 23.86\\

fixed & Capacity & 4.68 & 5.41 & 11.94 & 13.80\\

fixed & Rework X & 7.44 & 8.45 & 20.61 & 23.36\\

fixed & Rework Y & 7.46 & 7.99 & 20.59 & 22.05\\

service & None & 43.69 & 45.54 & 184.34 & 191.92\\

service & Capacity & 22.81 & 24.08 & 84.37 & 87.91\\

service & Rework X & 41.85 & 44.62 & 174.80 & 186.61\\

service & Rework Y & 42.17 & 42.49 & 174.97 & 174.70\\

routing & None & 31.13 & 30.37 & 111.82 & 109.49\\

routing & Capacity & 15.01 & 14.88 & 46.73 & 46.50\\

routing & Rework X & 16.54 & 21.86 & 54.25 & 72.12\\

routing & Rework Y & 21.54 & 16.27 & 70.68 & 52.47\\

combined & None & 91.64 & 89.79 & 567.45 & 554.47\\

combined & Capacity & 72.57 & 70.68 & 373.43 & 365.29\\

combined & Rework X & 69.66 & 83.69 & 355.70 & 471.85\\

combined & Rework Y & 86.18 & 69.81 & 491.10 & 366.32\\

\bottomrule\end{longtable}

\begin{longtable}{llrrr}\caption{P5: средние очереди B и D на исходном горизонте; число парных траекторий с перемещением наибольшей очереди от B к D.}\\\toprule

Mechanism & Action at 10 & Queue B & Queue D & B to D\\\midrule\endfirsthead

\toprule Mechanism & Action at 10 & Queue B & Queue D & B to D\\\midrule\endhead

fixed & None & 13.12 & 0.11 & 0/24\\

fixed & Capacity B & 1.99 & 4.75 & 24/24\\

fixed & Capacity D & 13.12 & 0.01 & 0/24\\

fixed & Rework & 13.12 & 0.08 & 0/24\\

service & None & 140.94 & 0.01 & 0/24\\

service & Capacity B & 38.36 & 0.09 & 0/24\\

service & Capacity D & 140.94 & 0.00 & 0/24\\

service & Rework & 140.94 & 0.00 & 0/24\\

routing & None & 13.28 & 35.73 & 0/24\\

routing & Capacity B & 2.09 & 46.05 & 0/24\\

routing & Capacity D & 13.28 & 4.24 & 0/24\\

routing & Rework & 13.28 & 1.69 & 0/24\\

combined & None & 141.02 & 1.49 & 0/24\\

combined & Capacity B & 40.19 & 118.56 & 24/24\\

combined & Capacity D & 141.02 & 0.05 & 0/24\\

combined & Rework & 141.02 & 0.03 & 0/24\\

\bottomrule\end{longtable}
